\chapter{Kết luận và hướng phát triển}
\label{ch:conclusion}

\section{Kết luận}

Trong bối cảnh thương mại điện tử ngày càng phát triển, đặc biệt ở nhóm ngành hàng thời trang
trẻ em, việc xây dựng một website bán hàng trực tuyến chuyên nghiệp là một hướng đi phù hợp,
mang lại lợi ích cho cả người bán lẫn người mua, đồng thời giảm thiểu các quy trình thủ công
truyền thống.

Xuất phát từ nhu cầu thực tiễn đó, đồ án đã thực hiện đề tài \textit{``\thesisTitle''}. Hệ thống được
thiết kế và xây dựng nhằm đáp ứng đầy đủ các chức năng của một website thương mại điện tử, bao
gồm: quản lý người dùng và phân quyền; quản lý sản phẩm theo danh mục, độ tuổi, giới tính và
biến thể; quản lý giỏ hàng, danh sách yêu thích và so sánh sản phẩm; đặt hàng và thanh toán trực
tuyến qua PayOS bên cạnh COD; quản lý đơn hàng, đổi/trả và tồn kho; chương trình khuyến mãi
(voucher, flash sale); đánh giá sản phẩm có kiểm duyệt; quản lý nội dung trang chủ; và thống kê
báo cáo doanh thu.

Về mặt kỹ thuật, hệ thống được xây dựng trên nền tảng công nghệ hiện đại (Next.js 14, React
Server Components, Prisma, PostgreSQL trên Supabase), có kiến trúc phân tầng rõ ràng, dễ bảo
trì và mở rộng, được kiểm thử tự động và đã sẵn sàng triển khai trên môi trường thực tế. Các
nghiệp vụ nhạy cảm như đặt hàng, trừ kho và áp dụng khuyến mãi được xử lý trong giao dịch nhằm
bảo đảm toàn vẹn dữ liệu; việc thanh toán được tích hợp an toàn với cơ chế tự động hủy đơn quá
hạn.

Kết quả đạt được của đồ án đã đáp ứng các mục tiêu và phạm vi nghiên cứu ban đầu. Thông qua
quá trình thực hiện, tôi đã củng cố kiến thức về phân tích và xây dựng hệ thống thông tin, đồng thời
rèn luyện nhiều kỹ năng quan trọng, bao gồm:

\begin{itemize}
  \item Thiết kế và xây dựng ứng dụng web full-stack theo kiến trúc nhiều tầng.
  \item Vận dụng quy trình phát triển phần mềm từ khảo sát, phân tích yêu cầu, thiết kế, cài đặt,
  kiểm thử đến triển khai.
  \item Thiết kế cơ sở dữ liệu quan hệ và xử lý các nghiệp vụ giao dịch phức tạp.
  \item Tích hợp dịch vụ bên ngoài (thanh toán, email) và bảo đảm an toàn ứng dụng.
  \item Kỹ năng nghiên cứu, tìm kiếm tài liệu và quản lý thời gian thực hiện dự án.
\end{itemize}

\section{Hướng phát triển}

Mặc dù đã đáp ứng các yêu cầu đặt ra ban đầu, do giới hạn về thời gian và nguồn lực, hệ thống vẫn
còn một số hướng có thể tiếp tục hoàn thiện và mở rộng trong tương lai:

\begin{itemize}
  \item \textbf{Đa dạng hóa phương thức thanh toán:} bổ sung thêm các cổng thanh toán phổ biến
  khác (ví điện tử, thẻ ngân hàng quốc tế) bên cạnh PayOS để tăng lựa chọn cho khách hàng.
  \item \textbf{Tích hợp đơn vị vận chuyển:} kết nối API của các đơn vị vận chuyển (GHN, GHTK)
  để tự động tính phí, tạo vận đơn và theo dõi hành trình giao hàng.
  \item \textbf{Phát triển ứng dụng di động:} xây dựng ứng dụng cho Android và iOS nhằm mở
  rộng khả năng tiếp cận và nâng cao trải nghiệm mua sắm.
  \item \textbf{Cá nhân hóa và gợi ý sản phẩm:} áp dụng các kỹ thuật gợi ý dựa trên lịch sử mua
  sắm và hành vi duyệt web để tăng tỷ lệ chuyển đổi.
  \item \textbf{Chăm sóc khách hàng:} bổ sung chat trực tuyến, thông báo đẩy (web push) và
  chương trình tích điểm, hạng thành viên để giữ chân khách hàng.
  \item \textbf{Phân tích dữ liệu nâng cao:} mở rộng hệ thống báo cáo với các chỉ số kinh doanh
  chuyên sâu và bảng điều khiển trực quan hỗ trợ ra quyết định.
\end{itemize}
